% ==============================================================================
% ANEXOS DA DISSERTAÇÃO (PROPECAUT / IFES)
% ==============================================================================

\newpage
\section*{ANEXOS}
\addcontentsline{toc}{section}{ANEXOS}

% ------------------------------------------------------------------------------
% ANEXO A
% ------------------------------------------------------------------------------
\subsection*{ANEXO A -- Especificações Técnicas de Qualificação da Célula Spectrolab XTJ Prime}
\addcontentsline{toc}{subsection}{ANEXO A -- Especificações Técnicas de Qualificação da Célula Spectrolab XTJ Prime}

Este anexo resume dados publicados pela Spectrolab para a célula solar de tripla junção XTJ Prime sob AM0 ($135{,}3\text{ mW/cm}^2=1353\text{ W/m}^2$, $28^\circ\text{C}$). O fabricante informa múltiplas geometrias; a área de $30{,}18\text{ cm}^2$ adotada no modelo não é apresentada aqui como dimensão universal de catálogo. Valores calculados a partir dessa área são identificados como tais.

\begin{table}[!ht]
\centering
\small
\renewcommand{\arraystretch}{0.85}
\caption{Extrato da folha de dados técnicos de qualificação espacial da célula Spectrolab XTJ Prime.}
\label{tab:anexo_datasheet_spectrolab}
\begin{tabular}{lcc}
\hline
\textbf{Característica Elétrica / Física (AM0, 28°C)} & \textbf{Valor Típico} & \textbf{Unidade} \\ \hline
Eficiência MPP BOL declarada & 30,7 & \% \\
Tensão de Circuito Aberto ($V_{oc}$) & 2,720 & V \\
Densidade de Corrente de Curto-Circuito ($J_{sc}$) & 18,0 & mA/cm$^2$ \\
Tensão no Ponto de Máxima Potência ($V_{mp}$) & 2,406 & V \\
Densidade de Corrente no MPP ($J_{mp}$) & 17,5 & mA/cm$^2$ \\
Coeficiente Térmico de Tensão de Circuito Aberto ($dV_{oc}/dT$) & $-5,60$ & mV/K \\
Coeficiente Térmico de Densidade de Corrente ($dJ_{sc}/dT$) & $+10,0$ & $\mu$A/(cm$^2$ K) \\
Coeficiente Térmico de Tensão MPP ($dV_{mp}/dT$) & $-6,30$ & mV/K \\
Coeficiente Térmico de Densidade de Corrente MPP ($dJ_{mp}/dT$) & $+5,0$ & $\mu$A/(cm$^2$ K) \\
Espessura Típica do Substrato de Germânio & 140 & $\mu$m \\
Densidade Mássica Superficial da Célula com Coverglass & 84 & mg/cm$^2$ \\
Tensão de Ruptura Reversa Típica ($V_{break}$) & $-25\text{ a }-30$ & V \\
\hline
\end{tabular}
\par\smallskip\noindent\textit{Nota:} Para a área adotada de $30{,}18\text{ cm}^2$, os dados $V_{mp}$ e $J_{mp}$ implicam $I_{mp}=0{,}52815\text{ A}$ e $P_{mp}=1{,}27073\text{ W}$ [CALCULADO]. Esses valores arredondados não reproduzem exatamente os $30{,}7\%$ de eficiência declarados no datasheet; a discrepância não foi recalibrada.
\end{table}

\vspace{0.8cm}
% ------------------------------------------------------------------------------
% ANEXO B
% ------------------------------------------------------------------------------
\subsection*{ANEXO B -- Diretrizes e Boas Práticas do Manual NASA-HDBK-4007A}
\addcontentsline{toc}{subsection}{ANEXO B -- Diretrizes e Boas Práticas do Manual NASA-HDBK-4007A}

O manual técnico de engenharia NASA-HDBK-4007A (\textit{Spacecraft High-Voltage Paschen and Corona Design Handbook}, 145 páginas, aprovado pela NASA em fevereiro de 2026) consolida as diretrizes de projeto e lições aprendidas para prevenção e mitigação de descargas por efeito corona, arcos de plasma no vácuo e quebra dielétrica em espaçonaves operando com barramentos de alta tensão ($V_{bus} > 100\text{ VDC}$). Embora não constitua uma norma de cumprimento mandatório genérico, o documento é adotado pelas missões da NASA e parceiros internacionais como o guia de referência para qualificação de subsistemas de potência HVDC.

\begin{table}[!ht]
\centering
\small
\renewcommand{\arraystretch}{0.85}
\caption{Diretrizes de mitigação e critérios de projeto do manual NASA-HDBK-4007A para barramentos de 400~V.}
\label{tab:anexo_nasa_hdbk_4007a}
\resizebox{0.95\textwidth}{!}{%
\begin{tabular}{p{5.5cm}cp{7.0cm}}
\hline
\textbf{Fenômeno / Requisito} & \textbf{Recomendação de Projeto} & \textbf{Diretriz de Mitigação no Satélite (ProPECAut/IFES)} \\ \hline
Tensão Crítica de Limiar de Paschen no Vácuo & $\ge 70\text{ a }100\text{ V}$ & Barramentos de 400~V operam acima do limiar mínimo de Paschen, exigindo isolamento sólido reforçado e supressão ativa de arco. \\
Controle de Energia de Arco e Danos & $E_{arc} < 5\text{ J}$ & Abertura ultrarrápida pelo interruptor SSPC em menos de $15\ \mu\text{s}$ para restringir a energia dissipada no curto-circuito. \\
Distâncias de Isolamento Dielétrico (Creepage/Clearance) & $> 0{,}8\text{ mm}$ a 400~V & Separação mínima entre trilhas e condutores não isolados expostos ao vácuo e plasma espacial na placa de circuito impresso. \\
Isolamento com Película de Kapton/Silicone & $\ge 2\text{ camadas}$ & Encapsulamento de chicotes, barramentos e conexões crimpadas para evitar descargas superficiais (\textit{flashover}). \\
Desgaseificação de Materiais (\textit{Outgassing}) & CVCM $< 0{,}1\%$ & Materiais isolantes e fitas térmicas devem atender a TML $< 1{,}0\%$ para evitar a formação de atmosfera local ionizável. \\
Proteção contra Corrente Reversa de Strings & Diodos / SSPC & Segmentação de strings com diodos de bloqueio e chaves estáticas de estado sólido para isolar células sombreadas. \\
\hline
\end{tabular}%
}
\end{table}
